\documentclass[10pt]{article}
\usepackage{booktabs}

\begin{document}

% Título del informe de laboratorio
\centering
Laboratorio de Química: Análisis de la Reacción entre $I^-$ y $(S_2O_4)^+$
\vspace*{10mm}

% Introducción
La reacción entre $I^-$ y $(S_2O_4)^+$ es una reacción interesante que se utiliza para estudiar las propiedades de los iones. En este informe, presentamos nuestros resultados sobre la velocidad de esta reacción.

% Tabla de resultados
\begin{table}[h!]
\centering    % centrar tabla
\begin{tabular}{@{}ccc@{}}
\toprule
Concentración de $I^-$ (M) & Concentración de $(S_2O_4)^+$ (M) & Velocidad de la reacción (M/s) \\ \midrule
0.1 & 0.1 & 0.01 \\ \bottomrule \\ 
\end{tabular}\\ \\ \\\textit{}
% Conclusiones
En conclusión, los resultados obtenidos en este informe muestran que la velocidad de la reacción entre $I^-$ y $(S_2O_4)^+$ es directamente proporcional a la concentración de los iones.

% Referencias

\end{document}
En conclusión, los resultados obtenidos en este informe muestran que la velocidad de la reacción entre $I^-$ y $(S_2O_4)^+$ es directamente proporcional a la concentración de los iones.

% Referencias

\end{document}